% Part: sets-functions-relations
% Chapter: functions
% Section: functions-relations

\documentclass[../../../include/open-logic-section]{subfiles}


\begin{document}

\olfileid[tr]{sfr}{fun}{rel}

\olsection{Bağıntı Olarak Fonksiyonlar}

\begin{explain}
$A$'nın elemanlarını $B$'nin elemanlarına eşleyen bir fonksiyon, açıktır ki
$A$ ile $B$ arasında bir bağıntı tanımlar: bu bağıntı $x$ ile $y$ arasında
ancak ve ancak $f(x)=y$ olduğunda geçerlidir. Hatta, örneğin matematiğin
yapıtaşlarını daha temel öğelere indirgemekle ilgileniyorsak, $f$
fonksiyonunu bu bağıntıyla, yani bir ikililer kümesiyle
\emph{özdeşleştirebiliriz}. Bu da şu soruyu ortaya çıkarır: hangi bağıntılar
bu yolla fonksiyon tanımlar?
\end{explain}

\begin{defn}[Bir fonksiyonun grafiği] $f\colon A \to B$ bir fonksiyon olsun.
$f$'nin \emph{grafiği},
\[
R_f = \Setabs{\tuple{x,y}}{f(x) = y}
\]
eşitliğiyle tanımlanan $R_f \subseteq A \times B$ bağıntısıdır.
\end{defn}

\begin{explain}
Kaplamsallık gereğince bir fonksiyonun grafiği biricik biçimde belirlenir.
Ayrıca kümeler için kaplamsallık, fonksiyonlara ilişkin örtük kaplamsallık
ilkesini hemen doğrular: $f$ ile $g$ aynı tanım ve değer kümelerine sahipse,
bütün değerlerde uyuştuklarında özdeştirler.

Benzer biçimde bir bağıntı ``fonksiyonel'' ise bir fonksiyonun grafiğidir.
\end{explain}

\begin{prop}\ollabel{prop:graph-function}
$R \subseteq A \times B$ bağıntısı aşağıdaki koşulları sağlasın:
\begin{enumerate}
\item $Rxy$ ve $Rxz$ ise $y=z$'dir; ve
\item her $x \in A$ için $\tuple{x,y}\in R$ olacak biçimde bir $y \in B$
vardır.
\end{enumerate}
O hâlde $R$, $f(x)=y$ ancak ve ancak $Rxy$ olduğunda geçerli olacak biçimde
tanımlanan $f\colon A \to B$ fonksiyonunun grafiğidir.
\end{prop}

\begin{proof}
İkinci koşul gereğince her $x \in A$ için $Rxy$ olacak biçimde en az bir
$y \in B$ vardır. Böyle bir $y$ seçelim. $Rxz$ olacak biçimde başka bir
$z\neq y$ bulunsaydı birinci koşul çiğnenirdi. Dolayısıyla her $x \in A$ için
$Rxy$ olacak biçimde bir ve yalnızca bir $y \in B$ vardır; böylece $f$ iyi
tanımlıdır. Açıktır ki $R_f=R$'dir.
\end{proof}

\begin{explain}
Her $f\colon A \to B$ fonksiyonunun bir grafiği, yani $f(x)=y$ eşitliğiyle
tanımlanan, $A$ ile $B$ arasında bir bağıntısı (bir $A\times B$ alt kümesi)
vardır. Öte yandan
\olref{prop:graph-function}'ta verilen özellikleri taşıyan her
$R \subseteq A \times B$ bağıntısı, bir $f \colon A \to B$ fonksiyonunun
grafiğidir. Fonksiyonlarla grafikleri arasındaki bu yakın bağlantı nedeniyle
bir fonksiyonu yalnızca grafiği olarak düşünebiliriz. Başka bir deyişle
fonksiyonlar, belirli bağıntılarla, yani belirli sıralı ikili kümeleriyle
özdeşleştirilebilir. \oliflabeldef{sfr:rel:ref:sec}{Yine de bu
``özdeşleştirmenin'' anlamı \olref[sfr][rel][ref]{sec}'teki gibidir: bu,
fonksiyonların metafiziği hakkında bir iddia değil, fonksiyonları belirli
kümeler olarak \emph{ele almanın} elverişli olduğuna ilişkin bir gözlemdir.
Bu özdeşleştirme elverişlidir; çünkü artık}{Artık} fonksiyonlar
üzerinde, bağıntılar üzerinde gerçekleştirdiğimize benzer işlemleri ele
alabiliriz (bkz. \olref[sfr][rel][ops]{sec}). Özellikle:
\end{explain}

\begin{defn}\ollabel{defn:funimage}
$f \colon A \to B$ bir fonksiyon ve $C\subseteq A$ olsun.

$f$'nin $C$'ye \emph{kısıtlaması}, her $x \in C$ için
$(\funrestrictionto{f}{C})(x)=f(x)$ eşitliğiyle tanımlanan
$\funrestrictionto{f}{C}\colon C \to B$ fonksiyonudur. Başka bir deyişle,
$\funrestrictionto{f}{C} = \Setabs{\tuple{x, y} \in R_f}{x \in C}$.

$f$'nin $C$'ye \emph{uygulanması}, $\funimage{f}{C} =
\Setabs{f(x)}{x \in C}$ kümesidir. Buna $C$'nin $f$ altındaki
\emph{görüntüsü} de denir.
\end{defn}

\begin{explain}
Bu tanımlardan, her $f$ fonksiyonu için $\ran{f} =
\funimage{f}{\dom{f}}$ eşitliği çıkar.
\oliflabeldef{sfr:rel:ops:sec}{Bu kavramlar, \olref[sfr][rel][ops]{sec}'teki
tanımlar ve fonksiyonları bağıntılarla özdeşleştirmemiz göz önüne alındığında,
tam da bekleneceği gibidir. Ancak diğer iki işlem, ters alma ve bağıntı
bileşkesi, biraz daha ayrıntı gerektirir. Bu ayrıntıları \olref[inv]{sec} ve
\olref[cmp]{sec}'te vereceğiz.}{}
\end{explain}

\end{document}
